\def\ChapterCode{EXPERT}
\chapter[Konventionelle und hydroaktive Wundauflagen]{Konventionelle und \\hydroaktive Wundauflagen}

 
\Z{
\noindent Hierbei werden nicht behandelt:
\begin{itemize}
  \item Antibakterielle Wundauflagen
  \item Die sehr saugstarken \I{Polyacrylat-Superabsorber}
  aus Acrylsäure und Natrium\-acrylat (z.\,B. \I{Sorbion
  sachet S{\texttrademark}})
  \item VAC-Systeme
  \item Interaktive/bioregulative Wundauflagen (Chitosan-/ Gelatine-haltige, Hyaluronsäure-haltige und Kollagen-Wundauflagen)
  \item Aktive Wundauflagen (Transplantate)
\end{itemize}
}





\section{Grundlagen}

\subsection{Wundheilung und Wundstadien}

Die normale Wundheilung verläuft in mehreren, überlappenden Stadien ab:

\begin{enumerate}
  \item \T{Reinigungsphase} (\I{Sekretionsphase}): Nach
  anfänglicher Blutstillung erfolgt eine \EE{Wundreinigung},
  welche mit einer gesteigerten Sekretion
  (Exsudation) einhergeht. Es bildet sich ein
  \EE{Fibrinnetz}, welches später als Matrix für die
  Granulation dient. In den ersten Stunden nach der Wundbildung beginnt eine
  \EE{Entzündungsreaktion} mit den typischen
  inflammatorischen Zeichen. Durch das Einwandern von
  Granulozyten und Monozyten beginnt die \EE{zelluläre
  Immunabwehr}. Im unkomplizierten Fall endet die
  Reinigungsphase am Tag 3.
  \item \T{Granulationphase}: Hier kommt es zum
  \EE{Gewebeaufbau} und zur Defektfüllung,
  \EE{Vaskularisation} und \EE{Granulation}. Die
  \EE{Fibroblasten} produzieren Kollagen und Proteoglykane,
  letztere bilden eine gallertartige Grundsubstanz für das Bindegewebe.
  Zink-abhängige Matrixmetallproteinasen fördern den Aufbau
  von Granulationsgewebe. \T{Granulationsgewebe} besteht vorwiegend aus
  Gewebsmakrophagen, Endothelzellen und Fibroblasten, welche
  von einer Matrix umgeben sind. Es hat eine \EE{körnige,
  feuchtglänzende Oberfläche} und ist in Folge der
  Vaskularisation \EE{dunkelrosa}.
  \item \T{Epithelisierungsphase}: Es kommt zu einem
  Ausreifen der Kollagenfasern und zur Umbildung des
  Granulationgewebes in Narbengewebe. Dabei kommt es durch
  \T{Wundkontraktion}  zu einer Verkleinerung der Wunde. Die
  Fibroblasen wandeln sich in Fibrozyten und Myofibroblasten
  um. Beim Prozess der Epithelisierung \EE{überhäuten vom
  Wundrand ausgehend Keratinozyten} mit einem
  \EE{feinen, rosafarbenen Epithelrasen} das
  Granulationsgewebe.
  Vorraussetzung dafür ist ein sauberes Granulationsgewebe,
  welches Wundrandniveau erreicht haben muss. Die
  Epithelisierungsphase ist nach etwa 3 Wochen beendet, die maximale Zugfestigkeit erhält
  die Wunde aber erst nach mehreren Wochen.
  \item \T{Remodellierungsphase}: Der endgültige Umbau des
  Narbengewebes kann Monate bis Jahre dauern.
\end{enumerate}




\section{Arten von Wundauflagen}

Ziel ist i.\,d.\,R. ein feuchtes Wundklima.
Die ideale Wundauflage erfüllt u.\,a. folgende Kriterien:

\begin{enumerate}
  \item Sie garantiert ein \EE{feuchtes Wundklima} (nicht
  gleichzusetzen mit Okklusion).
  \item Sie sorgt für die ausreichende \EE{Aufnahme von
  Wundsekret}.
  \item Sie ist \EE{nicht toxisch} und \EE{nicht allergen}.
  \item Sie verschmilzt nicht mit der Wunde, gibt keine
  Partikel ab.
  \item Sie bildet eine Barriere gegen Außeneinflüsse.
  \item Sie ist \EE{undurchlässig für Keime}.
  \item Sie ermöglicht einen \EE{Gasaustausch}.
\end{enumerate}

% \begin{description}
%   \item[Produkte (Beispiele)]
%   \item[Beschreibung] 
%   \item[Vorteile] ~
%   \begin{itemize}
%     \item[$+$] 
%   \end{itemize}
%   \item[Nachteile] ~
%   \begin{itemize}
%     \item[$-$] 
%   \end{itemize}
%   \item[Indikationen]~
%   \begin{itemize}
%     \item 
%   \end{itemize}
%   \item[Anwendung]
%   \item[Kontraindikationen] ~
%   \begin{itemize}
%     \item 
%   \end{itemize}
% \end{description}

\subsection{Koventionelle Wundauflagen}

\Text
{Mullkompressen}
{ 
\begin{description}
  \item[Beschreibung] Sie bestehen aus einem Baumwollgewebe,
  dessen Fadendichte je
nach Fabrikat variabel ist.
  \item[Vorteile] Sie sind billig und bieten eine
relativ gute Saugkraft. 
%   \begin{itemize}
%     \item[$+$] 
%   \end{itemize}
  \item[Nachteile] Als Nachteile sind jedoch Mazeration
infolge  flächiger Exsudataufnahme, die Notwendigkeit von
häufigen Verbandwechseln, die fehlende Keimbarriere und die
Gefahr der Austrocknung der Wunde zu nennen. Weiters
\EE{verkleben} sie mit dem Wundgrund, was einerseits den
Verbandswechsel schmerzhaft macht, und andererseits dabei
für ein neuerliches Wundtrauma sorgt.
%   \begin{itemize}
%     \item[$-$] 
%   \end{itemize}
  \item[Indikationen] Primärversorgung von Akutwunden, stark
  nässende Wunden mit häufigen Verbandswechsel,
  Arzneimittelträger; \EE{Sekundärverband}
  \item[Kontraindikationen] Primärverband bei allen sekundär
  heilenden, granulierenden, epithelisierenden oder schwach
  sezernierenden Wunden.
\end{description}
}
{}
{}
{}
{}

\Text
{Vliesstoffe}
{
\begin{description}
  \item[Beschreibung] Im Gegensatz zu Mullkompressen handelt
  es sich hierbei um \E{nicht-gewebte} Textilien. Es kommen
  unterschiedliche Fasern zum Einsatz.
  \item[Vorteile, Nachteile, Indikationen und Kontraindikationen] Siehe Mullkompressen.
\end{description}

}
{}
{}
{}
{}

\Text
{Saugkompressen}
{ 
\begin{description}
  \item[Beschreibung] Sie bestehen aus einem saugstarken
  Kern (z.\,B. aus Zellstoff-Flocken), welcher von einem
  Vliesstoff umhüllt ist.
  \item[Vorteile] Sie sind sind saugstark, weich und bieten eine gute Polsterung
%   \begin{itemize}
%     \item[$+$] 
%   \end{itemize}
  \item[Nachteile] Siehe Vliesstoffe. 
%   \begin{itemize}
%     \item[$-$] 
%   \end{itemize}
  \item[Indikationen] Primärversorgung von Akutwunden, stark
  nässender Wunden mit häufigen Verbandswechsel,
  Arzneimittelträger; Als \EE{Sekundärverband} gut geeignet
  bei stark sezernierenden Wunden, z.\,B. in Verbindung mit
  Alginaten.
  \item[Kontraindikationen] Mäßig bis schwach
  sezernierende Wunden, Primärverband bei granulierenden und epithelisierenden Wunden 
\end{description}
}
{}
{}
{}
{}


\subsection{Moderne Wundauflagen}

\subsubsection{Gaze und Wundfolien}


\TextSlide
{Imprägnierte Wundgaze}
{
\begin{description}
  \item[Produkte (Beispiele)] \I{Inadine{\texttrademark}}
  (Salbengaze mit PVP-Jod), \I{Jelonet{\texttrademark}}
  (Paraffingaze), \I{Mepitel{\texttrademark}}
  (Silikon-beschichtetes Polyamid-Netz)
  \item[Beschreibung] Imprägnierte Wundgaze sind
  grobmaschige Auflagen aus Natur- oder Kunstfasern, welche mit unterschiedlichen Stoffen 
(z.\,B. hydrophobe Substanzen (Paraffin, \ldots;
\enquote*{Fettgaze}), hydrophile Emulsionen, hydroaktive
bzw.
quellende Partikel, Silikon (\I{Mepithel{\texttrademark}}))
versetzt sind.
Die Imprägnierung verhindert das Verkleben der Wunde mit der dahinter
aufgebrachten Kompresse, durch die weiten Maschen ist dabei
eine ungehinderte Aufnahme von Exsudat gewährleistet.
  \item[Vorteile] ~
  \begin{itemize}
    \item[$+$] Preiswerte Alternative zu hydroaktiven
    Wundauflagen
    \item[$+$] Verklebungen minimiert
  \end{itemize}
  \item[Nachteile] ~
  \begin{itemize}
    \item[$-$] Fettgaze können bei geringen Exsudat
    anhaften, da sie hydrophob sind lassen sie sich auch
    nicht durch Spülung gut lösen
  \end{itemize}
  \item[Indikationen] Oberflächliche mäßig bis stark
  sezernierende Wunden
%   \begin{itemize}
%     \item 
%   \end{itemize}
  \item[Anwendung] Mind. täglicher Verbandwechsel
  \item[Kontraindikationen] Hydrophobe Gaze: Geringes Exsudat
%   \begin{itemize}
%     \item 
%   \end{itemize}
\end{description}
}
{
\begin{itemize}
  \item \I{Inadine{\texttrademark}}
  (Salbengaze mit PVP-Jod)
  \item \I{Jelonet{\texttrademark}}
  (Paraffingaze)
  \item \I{Mepitel{\texttrademark}}
  (Silikon-beschichtetes Polyamid-Netz)
\end{itemize}
}
{}
{}
{}

\subsubsection{}
\Text
{Semipermeable Wundfolien}
{
\begin{description}
%   \item[Produkte (Beispiele)]
  \item[Beschreibung] Es handelt sich hierbei um dünne,
  transparente Membranen aus Polyurethan. Sie sind
  durchlässig für Sauerstoff und Wasserdampf, bilden jedoch
  eine Barriere für Bakterien und Nässe. 
  \item[Vorteile] ~
  \begin{itemize}
    \item[$+$] Feuchtes Wundklima
    \item[$+$] Wundbeobachtung möglich
    \item[$+$] Wasserfest
  \end{itemize}
  \item[Nachteile] ~
  \begin{itemize}
    \item[$-$] Keine Saugkapazität
    \item[$-$] Ablösen kann problematisch sein
    \item[$-$] Keine Haftung auf feuchter Haut
  \end{itemize}
  \item[Indikationen]~
  \begin{itemize}
    \item Oberflächliche, nicht nässende Wunden,
    epithelisierende Wunden
    \item Fixation
    \item VAC-Versiegelung
  \end{itemize}
%   \item[Anwendung]
  \item[Kontraindikationen] Nässende oder infizierte Wunden
%   \begin{itemize}
%     \item
%   \end{itemize}
\end{description}
}
{}
{}
{}
{}

\subsubsection{Hydroaktive Wundauflagen}


\TextSlide
{Alginate}
{
\begin{description}
  \item[Produkte (Beispiele)] ~
\begin{itemize}
    \item Kalziumalginate: \I{Suprasorb A{\texttrademark}}
    \item Kalzium-Natriumalginate: \I{Seasorb
    Soft{\texttrademark}}
    \item CMC zugesetzt: \I{Seasorb Soft{\texttrademark}}
    \item Silberalginat: \I{Suprasorb A +
    Ag{\texttrademark}}
  \end{itemize}
  
  \item[Beschreibung] Alginate bestehen aus Kalzium-  bzw.
  Natriumsalzen der Alginsäure, welche aus Braunalgen
  gewonnen wird. 
  Die Kalzium-Alginatfasern nehmen Natrium-reiches Exsudat
  auf, bilden ein \EE{Gel} aus Natriumalginat und
  geben dabei \EE{Kalzium-Ionen} an die Umgebung ab. Das Gel
  ist hydrophil, bindet \EE{große Mengen an Flüssigkeit} und
  schließt Bakterien und Dedritus fest ein. Es gibt Produkte mit
  100\PC* Kalzium-Alginatanteil und oder mit einem
  Natrium-Alginat-Anteil von bis zu 20\PC*, was die
  Gelbildung beschleunigt.
  Die Dichte der Fasern ist ausschlaggebend für
  Gelierungsrate und für die Dauer der Saugfähigkeit. Durch
  Zusatz von \I{Carboxymethylcellulose} (\I{CMC}) wird ein
  sehr starkes und rasches Ansaugvermögen erreicht, es
  bildet sich ein formstabiles Gel. Die Ausbreitung von
  Exsudat und damit Mazerationen werden dadurch vermindert.
  \item[Vorteile] ~
  \begin{itemize}
    \item[$+$] Alginate sind sehr saugstark, sie können
    Sekret bis zum 20fachen des Eigengewichtes aufnehmen und
    in Gel umwandeln.
    \item[$+$] Feuchtes Wundklima
    \item[$+$] Gut anpassungsfähig, eignet sich zum (losen)
    tamponieren, quillt auf
    \item[$+$] Durch die Abgabe von Kalzium-Ionen wirken sie
    außerdem \EE{blutstillend}.
    \item[$+$] Gut geeignet für \EE{infizierte Wunden}
  \end{itemize}
  \item[Nachteile] ~
  \begin{itemize}
    \item[$-$] Ausreichende Exsudatmenge Voraussetzung,
    sonst Gefahr des Austrocknens
    \item[$~$] Gefahr der Mazeration, daher nicht über
    Wundraund hinaus
    \item[$-$] Sie müssen \EE{täglich erneuert} werden.
  \end{itemize}
  \item[Anwendung] \EE{Mit Sekundärverband}. Bei stark oder
  mäßig nässenden Wunden mit Saugkompressen, Verbandswechsel alle
  \Range{1}{3} Tage je nach Sekretion. Bei schwach
  sezernierenden Wunden Folienverband bzw. Hydrokolloid, das
  Alginat kann (theoretisch) bis zu 1 Woche in der Wunde
  verbleiben (aber CAVE Infektion!).
  \item[Indikationen] ~
\begin{itemize}
    \item Mäßig bis stark nässende Wunden
    \item Tiefe Wunden, Wundhöhlen und
    Wundtaschen, oberflächliche Wunden in der stark
    nässenden Reinigungsphase
    \item Infizierte und nicht infizierte chronische Wunden
    \item Blutende Wunden
    \item Combustio II\textdegree
  \end{itemize}
  \item[Kontraindikationen] Trockene und nekrotische Wunden,
  Combustio III\textdegree
\end{description}
}
{
\begin{itemize}
    \item Kalziumalginate: \I{Suprasorb A{\texttrademark}}
    \item Kalzium-Natriumalginate: \I{Seasorb
    Soft{\texttrademark}}
    \item CMC zugesetzt: \I{Seasorb Soft{\texttrademark}}
    \item Silberalginat: \I{Suprasorb A +
    Ag{\texttrademark}}
  \end{itemize}
}
{}
{}
{}

\TextSlide
{Hydrokolloide}
{
\begin{description}
  \item[Produkte (Beispiele)] \I{Suprasorb
  H{\texttrademark}},
  \I{Varihesive{\texttrademark}}
  \item[Beschreibung] Hydrokolloide bestehen aus
  Natriumcarboxymethyzellulose, Gelatine, Pektin, Elastomeren und Klebstoffen und
  leigen in unterschiedlichen Varianten vor, z.\,B. als
  Platten oder Paste.
  Sie verwandeln \EE{Exsudat in ein Gel}, welches
  trübe und übel riechend sein kann. Sie sind undurchlässig für
  Dampf und Luft und fördern daher die \EE{Rehydrierung} von
  trockenen Wunden, sie zählen zu den Okklusivverbänden. Durch
  die vorübergehende lokale Hypoxie werden Makrophagen und die
  Angiogenese angeregt. 
  \item[Vorteile] ~
  \begin{itemize}
    \item[$+$] Feuchtigkeitsabgabe
    \item[$+$] Unterstützung des autolytischen Débridements
    \item[$+$] Selbstklebend, kein Sekundärverbad notwendig
    \item[$+$] Verweildauer bis zu 7 Tage
  \end{itemize}
  \item[Nachteile] ~
  \begin{itemize}
    \item[$-$] Irritationen und Allergien
    \item[$-$] Starke Haftung auf trockener Haut (Cave
    empfindliche Haut!), schlechte Haftung auf feuchter Haut
  \end{itemize}
  \item[Indikationen] ~
  \begin{itemize}
    \item Leicht bis stark nässende Wunden, je nach
    Schichtdicke
    \item \EE{Auflösen} von fibrinösen und schmierigen
    Belegen
    \item Einsetzbar in allen Wundheilungsphasen
    \item \T{Hydrokolloidplatten} eignen sich gut
  für Wundhöhlen oder flache Wunden mit mäßiger Sekretion und zeichnen  sich durch
  eine gute Anpassbarkeit aus. \T{Hydrokolloidpaste}
  kann einige Tage in der Wunde belassen werden und wirkt
  granulationsfördernd. 
  \end{itemize}
  \item[Kontraindikationen] Infizierte Wunden, Ischämische
  Ulzera, Wunden mit freiliegenden Muskeln, Sehnen oder
  Knochen
\end{description}
}
{
\begin{itemize}
  \item \I{Suprasorb
  H{\texttrademark}}
  \item 
  \I{Varihesive{\texttrademark}}
\end{itemize}
}
{}{}{}

\TextSlide
{Hydrofaser}
{
\begin{description}
  \item[Produkte (Beispiele)] \I{Aquacel{\texttrademark}} 
  \item[Beschreibung] {Hydrofasern} sind textile
  Gespinnste von Hydrokolloiden, welche bei der Aufnahme von
  Flüssigkeit ein \EE{formstabiles Gel} bilden
  und aufquellen.
  Sie weisen eine
  \EE{starke Absorptionsfähigkeit} auf und eignen sich für
  flache und auch ausgehöhlte Wunden mit mäßig bis starker Sekretion und
  können einige Tage in der Wunde belassen werden.
  \item[Vorteile] ~
  \begin{itemize}
    \item[$+$] Große Saug- und
    Speicherkapazität
    \item[$+$] Verweildauer theoretisch bis zu 7 Tage
  \end{itemize}
  \item[Nachteile] Exsudat notwendig, sonst Gefahr des
    Austrocknens
%   \begin{itemize}
%     \item[$-$] 
%   \end{itemize}
  \item[Indikationen] Mäßig bis stark nässende 
    Wunden
%   \begin{itemize}
%     \item 
%   \end{itemize}
  \item[Kontraindikationen] Trockene oder wenig
  sezernierende Wunden
\end{description}
}
{
\begin{itemize}
  \item \I{Aquacel{\texttrademark}} 
\end{itemize}
}{}{}{}

\TextSlide
{Hydrogele}
{
\begin{description}
  \item[Produkte (Beispiele)] Verbände: \I{Suprasorb
  G{\texttrademark}}; Tuben: \I{Suprasorb G
  Gel{\texttrademark}}, \I{Varihesive
  Hydrogel{\texttrademark}}
  \item[Beschreibung] Hydrogele bestehen aus einer Matrix
  unlöslicher hydrophiler Polymere, welche einen \EE{sehr
  hohen Wasseranteil} (bis zu 96\PC*) aufweisen. Sie sind
  grundsätzlich durchlässig für Wasserdampf und Sauerstoff.
  Sie können durch \E{langsame}, aber kontinuierliche Quellung
  auch (begrenzt) Sekret aufnehmen.
  
  Ihre Stärke ist das Feucht\E{halten} von schwach
  exsudierenden Wunden und das \EE{Aufweichen von Schorf und
  Belegen}, sowie als Unterstützung des autolytischen
  Débridements.
  Sie sind als Gel oder Gelkompressen erhältlich. 
  \item[Vorteile] ~
  \begin{itemize}
    \item[$+$] Feuchthalten der Wunde
    \item[$+$] Autolyse
    \item[$+$] Kühlender Effekt bei Combustio
    \Range{I}{II}\textdegree
    \item[$+$] Gelkompressen sind durchsichtig
    (\enquote*{Schaufenster})
  \end{itemize}
  \item[Nachteile] ~
  \begin{itemize}
    \item[$-$] Eingeschränkte Suagfähigkeit,
    Mazerationsgefahr
    \item[$-$] Bei pAVK-Ulzera evtl. schmerzhaft
  \end{itemize}
  \item[Indikationen]~
  \begin{itemize}
    \item Gut bei schmierigen und nekrotischen Wunden
    \item Aufweichen von Nekrosen und Belegen
    \item Schwach bis mäßig nässende Wunden in Granulations-
    und Epithelisierungsphase
  \end{itemize}
  \item[Anwendung] Kompressen: Verbandwechsel je nach
  Exsudat alle \Range{1}{7} Tage
  
  Gel: Verbandswechselintervall bis zu 3 Tage 
  \item[Kontraindikationen] ~
  \begin{itemize}
    \item Starke Sekretion, Blutungen
    \item Cave bei infizierten Wunden
  \end{itemize}
\end{description}
}
{
\begin{itemize}
  \item Verbände:
  \begin{itemize}
    \item \I{Suprasorb
    G{\texttrademark}}
  \end{itemize}
  \item Tuben:
  \begin{itemize}
    \item \I{Suprasorb G
    Gel{\texttrademark}}
    \item \I{Varihesive
    Hydrogel{\texttrademark}}
  \end{itemize}
\end{itemize}
}
{}
{}
{}





\TextSlide
{Schaumstoffe und Hydropolymere}
{
\begin{description}
  \item[Produkte (Beispiele)] \I{Allevyn{\texttrademark}},
  \I{Mepilex{\texttrademark}}, \I{Suprasorb
  P{\texttrademark}}
  \item[Beschreibung] Schaumstoffe auf \I{Polyurethan}-Basis
  (\I{PU}) sind dank ihres \EE{Porenreichtums saugstarke,
  biologisch reizlose} Wundauflagen. Reine PU-Kompressen
  behalten dabei ihre Form bei, wohingegen Hydropolymere bei
  Exsudataufnahme expandieren. \Z{(Anmerkung: offenporige
  Schaumstoffkompressen werden an dieser Stelle nicht
  besprochen.)}
  \item[Vorteile] ~
  \begin{itemize}
    \item[$+$] Aufnahmekapazität bis zum
    \EE{\Range{20}{30}-fachen} des Eigengewichts
    \item[$+$] Thermoisolation, Polsterung
    \item[$+$] Feuchtes Wundklima, freier
    Gas-/ Wasserdampfaustausch
    \item[$+$] Kein Anhaften
  \end{itemize}
  \item[Nachteile] ~
  \begin{itemize}
    \item[$\sim$] Exsudat notwendig
    \item[$-$] Mazerationsgefahr
    \item[$-$] Zähflüssiges Sekret problematisch (Poren
    verstopft)
  \end{itemize}
  \item[Indikationen]~
  \begin{itemize}
    \item Mäßig bis stark sezernierende Wunden
  \end{itemize}
%   \item[Anwendung]
  \item[Kontraindikationen] ~
  \begin{itemize}
    \item Trockene oder nekrotische Wunden
    \item Tiefe Wunden mit freiliegenden Knochen, Sehnen
    oder Muskeln
  \end{itemize}
\end{description}
}
{
\begin{itemize}
  \item \I{Allevyn{\texttrademark}}
  \item \I{Mepilex{\texttrademark}} 
  \item \I{Suprasorb
  P{\texttrademark}}
\end{itemize}
}
{}
{}
{}

\subsection{Exkurs: Auf das Detail kommt es an}

Die Geschichte mancher Wundauflagen ist eine Geschichte
voller Missverständnisse:

 
\begin{table}[H]
\caption{Suprasorb{\texttrademark}: Namensschema mit
Verwechslungspotential }


\begin{tabular}{ll}
\addlinespace\toprule\addlinespace
\noindent\TabHeading{Handelsname} & \noindent\TabHeading{Art}
\\\addlinespace\midrule\addlinespace
\I{Suprasorb A{\texttrademark}}			& Alginat
\\
\I{Suprasorb A + Ag{\texttrademark}}	& Silberhaltige Auflage
\\
\I{Suprasorb C{\texttrademark}}			& Kollagen-Wundauflage
\\
\I{Suprasorb F{\texttrademark}}			& Semipermeable Wundfolie
\\
\I{Suprasorb G{\texttrademark}}			& Hydrogel 
\\
\I{Suprasorb G Gel{\texttrademark}}	& Hydrogel (Tube)
\\
\I{Suprasorb M{\texttrademark}}			& Sonstiges
\\
\I{Suprasorb P{\texttrademark}}			&	Schaumstoff
\\
\I{Suprasorb X{\texttrademark}}			& Sonstiges
\\
\I{Suprasorb X + PHMB{\texttrademark}}		& Sonstiges
\\\addlinespace\bottomrule\addlinespace
\end{tabular}


\end{table}




\section{Auswahl der richtigen Wundauflage}

\Layout{57}
{}
{}
{}
{\begin{tabularx}{\linewidth}{XXXXXXX}
\toprule\addlinespace
\noindent\TabHeading{Typ}
&
\multicolumn{4}{c}{\noindent\TabHeading{Reinigung}}
&
\noindent\TabHeading{Granulation}
&
\noindent\TabHeading{Epithelisierung}
\\%\addlinespace
&
\noindent\TabHeading{Blutend}
&
\noindent\TabHeading{Exsudativ}
&
\noindent\TabHeading{Belegt}
&
\noindent\TabHeading{Infiziert}
&
&
\\\addlinespace\midrule\addlinespace
\noindent\TabSub{Saugkompressen}		
&	$+$		& $++$		& $\sim$	& $\sim$	& 				&		
\\\addlinespace
\noindent\TabSub{}		
&				& 				& 				& 				& 				&		
\\\addlinespace
\noindent\TabSub{Impräg. Gaze}		
&	$+$			& 	$+$			& 				& 	& 	$\sim$			&		
\\\addlinespace
\noindent\TabSub{Kohlekompressen}		
&				& 				& 				& $++$		& 				&		
\\\addlinespace
\noindent\TabSub{Alginate}		
& $++$		& $++$		& $+$		& $++$		& $+$		&		
\\\addlinespace
\noindent\TabSub{Hydrofiber}		
&	$+$		& $++$		& 				& 				& 	$+$		&		
\\\addlinespace
\noindent\TabSub{Hydrokolloide}		
&				& 	$+$		& $+$		& 				& 	$++$	&	$+$	
\\\addlinespace
\noindent\TabSub{Hydrogele}		
&				& 				& 	$++$	& 				& $++$		&	$+$	
\\\addlinespace
\noindent\TabSub{Schaumstoff}		
&				& 	$++$	& $+$		& $\sim$	& $+$		&		
\\\addlinespace
\TabSub{Semipermeable Wundfolien}		
&				& 				& 				& 				& 				&	$++$	
\\\addlinespace
\noindent\TabSub{Silberhaltige Wundauflagen}		
&			& 			& 			& $++$		& 			&		
% \\\addlinespace
\\\addlinespace\bottomrule\addlinespace
\end{tabularx}}
{Phasengerechter Einsatz von Wundauflagen}
{nach 
\cite{AerztlicheWundversorgung2}
}


\Layout{57}
{}
{}
{}
{\begin{tabularx}{\linewidth}{XXXX}
\toprule\addlinespace
\noindent\TabHeading{Exsudatmenge}
&
\noindent\TabHeading{Flache Wunde}
&
\noindent\TabHeading{Tiefe Wunde}
&
\noindent\TabHeading{Wundhöhle}
\\\addlinespace\midrule\addlinespace
\TabSub{Blutend}
&
\begin{compactitem}
  \item Alginat
  \item Gelatineschwamm
  \item Kollagenschwamm
\end{compactitem}
&
\begin{compactitem}
  \item Alginat
  \item Gelatineschwamm
  \item Kollagenschwamm
\end{compactitem}
&
\begin{compactitem}
  \item Alginattamponade
  \item Gelatineschwamm
  \item Kollagenschwamm
\end{compactitem}
\\\addlinespace
\TabSub{Extrem nässend}
&
\begin{compactitem}
  \item Saugkompresse
  \item Superabsorber
\end{compactitem}
&
\begin{compactitem}
  \item Saugkompresse
  \item Superabsorber
\end{compactitem}
&
\begin{compactitem}
  \item Mullkompressen
  \item Superabsorber
\end{compactitem}
\\\addlinespace
\TabSub{Stark nässend}
&
\begin{compactitem}
  \item Alginat
  \item Hydrofiber
  \item Hydropolymer
  \item Schaumstoffkompresse
  \item Superabsorber
\end{compactitem}
&
\begin{compactitem}
  \item Alginat
  \item Hydrofiber
  \item \enquote*{Cavity}-Formen von Hydropolymer oder
  Schaumstoffkompresse
  \item Superabsorber
\end{compactitem}
&
\begin{compactitem}
  \item Alginattamponade
  \item Hydrofiber
  \item \enquote*{Cavity}-Formen von Hydropolymer oder
  Schaumstoffkompresse
  \item Superabsorber
\end{compactitem}
\\\addlinespace
\TabSub{Mäßig}
&
\begin{compactitem}
  \item Hydrokolloid
  \item Alginat
  \item Hydrofiber
  \item Hydropolymer
  \item Schaumstoffkompresse
  \item Nasstherapeutikum
\end{compactitem}
&
\begin{compactitem}
  \item Alginat
  \item Hydrofiber
  \item \enquote*{Cavity}-Formen von Hydropolymer oder
  Schaumstoffkompresse
  \item Nasstherapeutikum
\end{compactitem}
&
\begin{compactitem}
  \item Alginattamponade
  \item Hydrofiber
  \item \enquote*{Cavity}-Formen von Hydropolymer oder
  Schaumstoffkompresse
  \item Nasstherapeutikum
\end{compactitem}
\\\addlinespace
\TabSub{Wenig}
&
\begin{compactitem}
  \item Hydrogelkompresse
  \item Hydrokolloid dünn
  \item Semipermeable Wundfolie
\end{compactitem}
&
\begin{compactitem}
  \item Hydrogel
  \item Alginat (angefeuchtet)
\end{compactitem}
&
\begin{compactitem}
  \item Hydrogel
  \item Alginattamponade (angefeuchtet)
\end{compactitem}

\\\addlinespace\bottomrule\addlinespace
\end{tabularx}

% \end{WidePar}
}
{Exsudatmenge vs. Wundtiefe}
{nach 
\cite{AerztlicheWundversorgung2}
}
% \begin{table}
% \caption{}
% % \begin{WidePar}
% 
% \end{table}



\begin{Literary}
\EE{Exzellentes und umfassendes
Werk zur Wundversorgung:} 
\fullcite{WundversorgungPflege2}

\EE{Prägnante, weniger detaillierte Darstellung:} 
\fullcite{AerztlicheWundversorgung2}

\minisec{Weiters}

\fullcite{ThiemesPflege12}

\fullcite{SiewertChirurgie8}

\fullcite{GeschichteNeuerenDeutschenChirurgie}
\end{Literary}


\ChapterEnd
